\documentclass[12pt,a4paper]{article}
\usepackage[T1]{fontenc}
\usepackage[utf8]{inputenc}
\usepackage[brazil]{babel}
\usepackage{amsmath,amssymb,amsfonts,amsthm}
\usepackage{geometry}
\usepackage{graphicx}
\usepackage{booktabs}
\usepackage{array}
\usepackage{fancyhdr}
\usepackage{titlesec}
\usepackage{xcolor}
\usepackage{setspace}
\usepackage{mathtools}
\usepackage{bm}
\usepackage{hyperref}
\usepackage{enumitem}
\usepackage{tcolorbox}
\usepackage{mdframed}
\usepackage{listings}

\geometry{a4paper, top=3cm, bottom=2.5cm, left=3cm, right=2cm, headheight=25pt}
\setstretch{1.35}

\definecolor{ufamblue}{RGB}{0,51,102}
\definecolor{ufamgreen}{RGB}{0,102,51}
\definecolor{resultbg}{RGB}{240,248,255}
\definecolor{lightgray}{RGB}{245,245,245}
\definecolor{codebg}{RGB}{248,248,248}
\definecolor{codekw}{RGB}{0,51,153}
\definecolor{codecmt}{RGB}{110,110,110}
\definecolor{codestr}{RGB}{140,10,10}

\hypersetup{colorlinks=true, linkcolor=ufamblue, urlcolor=ufamblue, citecolor=ufamblue}

\pagestyle{fancy}
\fancyhf{}
\fancyhead[L]{\small\textcolor{ufamblue}{\textbf{PPGEE/UFAM} --- Verif.\ e S\'intese de Sist.\ Ciber-F\'isicos}}
\fancyhead[R]{\small\textcolor{ufamblue}{Lista 2 --- DSVerifier \& $k$-Indu\c{c}\~ao}}
\fancyfoot[C]{\small\thepage}
\renewcommand{\headrulewidth}{0.6pt}
\renewcommand{\headrule}{\hbox to\headwidth{\color{ufamblue}\leaders\hrule height \headrulewidth\hfill}}

\titleformat{\section}[block]
  {\large\bfseries\color{ufamblue}}
  {Exerc\'icio \thesection.}{0.8em}{}
  [\vspace{-0.3em}\textcolor{ufamblue}{\rule{\linewidth}{0.6pt}}]
\titlespacing*{\section}{0pt}{2.5ex}{1.5ex}

\titleformat{\subsection}{\normalsize\bfseries\color{ufamgreen}}{\thesubsection}{0.8em}{}
\titlespacing*{\subsection}{0pt}{1.5ex}{0.8ex}

\tcbset{
  resultbox/.style={
    colback=resultbg, colframe=ufamblue, fonttitle=\bfseries\small,
    title=Resultado, boxrule=0.8pt, arc=3pt,
    left=6pt, right=6pt, top=4pt, bottom=4pt
  },
  notabox/.style={
    colback=orange!5, colframe=orange!60!black, fonttitle=\bfseries\small,
    title=Nota did\'atica, boxrule=0.7pt, arc=3pt,
    left=6pt, right=6pt, top=4pt, bottom=4pt
  },
  quotebox/.style={
    colback=gray!6, colframe=gray!55, fonttitle=\bfseries\small,
    title=Cita\c{c}\~ao direta, boxrule=0.6pt, arc=2pt,
    left=6pt, right=6pt, top=3pt, bottom=3pt
  }
}

\lstdefinestyle{cstyle}{
  language=C, backgroundcolor=\color{codebg}, basicstyle=\ttfamily\footnotesize,
  keywordstyle=\color{codekw}\bfseries, commentstyle=\color{codecmt}\itshape,
  stringstyle=\color{codestr}, numbers=left, numberstyle=\tiny\color{gray},
  numbersep=6pt, frame=single, rulecolor=\color{gray!40}, breaklines=true,
  showstringspaces=false, tabsize=2, xleftmargin=1em
}
\lstdefinestyle{shstyle}{
  backgroundcolor=\color{codebg}, basicstyle=\ttfamily\footnotesize,
  frame=single, rulecolor=\color{gray!40}, breaklines=true,
  showstringspaces=false, xleftmargin=1em
}

\newcommand{\vv}[1]{\texttt{#1}}

\begin{document}

%==============================================================
% CAPA
%==============================================================
\begin{titlepage}
\centering
\vspace*{0.5cm}

{\color{ufamblue}\rule{\linewidth}{2pt}}\\[0.4cm]
{\Large\bfseries\color{ufamblue} UNIVERSIDADE FEDERAL DO AMAZONAS}\\[0.2cm]
{\large\color{ufamblue} Programa de P\'os-Gradua\c{c}\~ao em Engenharia El\'etrica --- PPGEE}\\[0.1cm]
{\large\color{ufamblue} Linha 1 --- Sistemas Inteligentes e Microeletr\^onica}\\[0.4cm]
{\color{ufamblue}\rule{\linewidth}{2pt}}\\[2.5cm]

{\LARGE\bfseries Lista de Exerc\'icios 02}\\[0.4cm]
{\Large\bfseries Verifica\c{c}\~ao de Sistemas Digitais e $k$-Indu\c{c}\~ao}\\[0.4cm]
{\large PGENE601/PPGINF554 --- Verifica\c{c}\~ao e S\'intese Autom\'atica}\\[0.1cm]
{\large de Sistemas Ciber-F\'isicos e Embarcados}\\[0.4cm]
{\normalsize Semana 5 (28/09/2026 a 04/10/2026)}\\[3cm]

\begin{tabular}{rl}
  \textbf{Disciplina:} & Verifica\c{c}\~ao e S\'intese Autom\'atica de Sist.\ Ciber-F\'isicos e Emb. \\[0.3cm]
  \textbf{Docente:}    & Prof.\ Dr.\ Lucas C.\ Cordeiro \\[0.3cm]
  \textbf{Discente:}   & Matheus Serr\~ao Uch\^oa \\[0.3cm]
  \textbf{N\'ivel:}    & P\'os-Gradua\c{c}\~ao --- Mestrado \\[0.3cm]
  \textbf{Artigos-base:} & DSVerifier (SPIN 2015) e $k$-Indu\c{c}\~ao (STTT 2017) \\[0.3cm]
  \textbf{Data:}       & Manaus, \today \\
\end{tabular}

\vfill
{\color{ufamblue}\rule{\linewidth}{1.5pt}}
\end{titlepage}

\newpage
\tableofcontents
\newpage

%==============================================================
% ARTIGOS DE REFERÊNCIA
%==============================================================
\section*{Artigos de Refer\^encia}
\addcontentsline{toc}{section}{Artigos de Refer\^encia}

\textbf{Artigo A.} H.\,I.\ Ismail, I.\,V.\ Bessa, L.\,C.\ Cordeiro, E.\,B.\ de Lima Filho e
J.\,E.\ Chaves Filho. \emph{DSVerifier: A Bounded Model Checking Tool for Digital Systems}.
SPIN 2015, LNCS 9232, pp.~126--131.

\textbf{Artigo B.} M.\,Y.\,R.\ Gadelha, H.\,I.\ Ismail e L.\,C.\ Cordeiro. \emph{Handling Loops
in Bounded Model Checking of C Programs via $k$-Induction}. STTT, 2017.

Todas as respostas abaixo citam diretamente o texto dos artigos (caixas cinzas) sempre que a
justificativa depende da formula\c{c}\~ao exata usada pelos autores.

\newpage

%==============================================================
% EXERCÍCIO 1
%==============================================================
\section{Efeitos de \emph{finite word length} (Artigo A)}

\subsection*{Enunciado}
Controlador de segunda ordem $H(z) = \dfrac{2{,}813z^2 - 0{,}0163z - 1{,}872}{z^2 + 1{,}068z + 0{,}1239}$,
planta de motor de indu\c{c}\~ao, per\'iodo de amostragem $0{,}5\,$s, formato \vv{int\_bits=4},
\vv{frac\_bits=10}, \vv{min=-5.0}, \vv{max=5.0}.

\subsection{(a) Faixa representável e resolução do formato $\langle 4,10\rangle$}

O Artigo A usa a conven\c{c}\~ao de ponto fixo em que \vv{int\_bits} \emph{inclui} o bit de
sinal (a struct \vv{implementation} declara diretamente \vv{int\_bits} e \vv{frac\_bits}, sem
um campo separado de sinal). Em complemento de dois, com $I=4$ bits inteiros (sinal incluso) e
$F=10$ bits fracion\'arios, a palavra tem $I+F=14$ bits e representa valores
$v = m \cdot 2^{-F}$, onde $m$ \'e um inteiro em complemento de dois de $I+F$ bits:
\[
  m \in \bigl[-2^{I+F-1},\; 2^{I+F-1}-1\bigr] = [-8192,\; 8191].
\]

\begin{align*}
  v_{\min} &= -2^{I-1} = -2^{3} = -8{,}000000 \\
  v_{\max} &= 2^{I-1} - 2^{-F} = 8 - 2^{-10} = 8 - 0{,}0009765625 = 7{,}9990234375 \\
  \Delta   &= 2^{-F} = 2^{-10} = 0{,}0009765625 \approx 9{,}7656\times10^{-4}
\end{align*}

\begin{tcolorbox}[resultbox]
Faixa representável $\langle4,10\rangle$: $[-8{,}0;\ 7{,}999023]$. Resolu\c{c}\~ao (LSB):
$\Delta = 2^{-10} \approx 0{,}000977$.
\end{tcolorbox}

\subsection{(b) Os coeficientes e a faixa $[-5,5]$ cabem em $\langle4,10\rangle$?}

Os coeficientes de $H(z)$, conforme o pr\'oprio arquivo de especifica\c{c}\~ao do DSVerifier
(\vv{ds.b = \{2.813, -0.0163, -1.872\}}, \vv{ds.a = \{1.0, 1.068, 0.1239\}}), e a faixa de
entrada $[-5{,}0;\,5{,}0]$:

\begin{center}
\begin{tabular}{@{}lcc@{}}
\toprule
Valor & Dentro de $[-8{,}0;\,7{,}999023]$? & \\
\midrule
$b_0=2{,}813$   & Sim & $-8 \le 2{,}813 \le 7{,}999$ \\
$b_1=-0{,}0163$ & Sim & \\
$b_2=-1{,}872$  & Sim & \\
$a_0=1{,}0$     & Sim & \\
$a_1=1{,}068$   & Sim & \\
$a_2=0{,}1239$  & Sim & \\
$[-5{,}0;\,5{,}0]$ & Sim & $[-5,5]\subset[-8;7{,}999]$ \\
\bottomrule
\end{tabular}
\end{center}

\begin{tcolorbox}[resultbox]
\textbf{Tudo cabe.} Nenhum coeficiente nem a faixa de entrada excede $[-8{,}0;\,7{,}999023]$,
logo $\langle4,10\rangle$ \'e um formato v\'alido em termos de \emph{faixa} (n\~ao h\'a
overflow estrutural nos valores nominais). Isso n\~ao implica representa\c{c}\~ao
\textbf{exata}: por exemplo $0{,}1239$ n\~ao \'e um m\'ultiplo de $2^{-10}$ e ser\'a
arredondado para o m\'ultiplo mais pr\'oximo, $130/1024 \approx 0{,}126953$ (erro de
quantiza\c{c}\~ao $\approx 3\times10^{-3}$) --- esse \'e precisamente o efeito de
\emph{finite word length} que o Artigo A se prop\~oe a verificar (overflow \emph{durante a
execu\c{c}\~ao} do filtro, quando somas/produtos intermedi\'arios de estados internos, n\~ao
apenas os coeficientes nominais, podem extrapolar a faixa).
\end{tcolorbox}

\subsection{(c) Por que a inicialização precisa ocorrer antes da instrumentação}

\begin{tcolorbox}[quotebox]
``DSVerifier performs three main procedures: initialization, validation, and instrumentation
[\ldots] the first step is to initialize its internal parameters for quantization, that is, it
computes the maximum and minimum representable numbers for the chosen FWL format. [\ldots] In
the last step, DSVerifier adds explicit calls to the verification engine (for the evaluated
properties), using functions available in ESBMC (e.g., \_\_ESBMC\_assume and
\_\_ESBMC\_assert), in order to check for property violations.'' (Artigo A, Se\c{c}\~ao~2)
\end{tcolorbox}

A instrumenta\c{c}\~ao \'e o passo que \textbf{escreve, como texto C literal}, as chamadas
\vv{\_\_ESBMC\_assert(x <= max \&\& x >= min)} (e equivalentes para overflow/limit
cycle) dentro do arquivo ANSI-C gerado. Para emitir essas linhas, a instrumenta\c{c}\~ao
precisa j\'a conhecer os valores num\'ericos concretos de \vv{max} e \vv{min} --- eles
entram como \textbf{constantes literais} no c\'odigo gerado, n\~ao como express\~oes a
resolver depois.

\begin{tcolorbox}[notabox]
O ESBMC, que roda \textbf{depois} (Fig.~1 do Artigo A: C Parser $\to$ GOTO Program $\to$ GOTO
Symex $\to$ SMT Solver), \'e um \emph{model checker de C gen\'erico}: ele n\~ao tem nenhuma
no\c{c}\~ao embutida de ``ponto fixo'', ``FWL'' ou ``formato $\langle k,l\rangle$'' --- s\'o
enxerga \vv{assert}/\vv{assume} com constantes literais. Se o c\'alculo de \vv{max}/\vv{min}
fosse adiado para ``durante a verifica\c{c}\~ao'', n\~ao haveria mais nenhum componente no
pipeline capaz de calcul\'a-lo: o DSVerifier j\'a delegou o controle ao ESBMC nesse ponto, e o
arquivo ANSI-C j\'a foi escrito e fechado. Por isso a inicializa\c{c}\~ao \'e necessariamente
uma etapa de \emph{gera\c{c}\~ao de c\'odigo} (tempo de compila\c{c}\~ao/instrumenta\c{c}\~ao),
estritamente anterior e externa \`a verifica\c{c}\~ao propriamente dita.
\end{tcolorbox}

\subsection{(d) Oscilação persistente com entrada constante}

\begin{tcolorbox}[quotebox]
``Limit Cycle. There can be persistent oscillations in the output of a system with constant
input. DSVerifier is able to check for zero-input limit cycles, for any initial condition.''
(Artigo A, Se\c{c}\~ao~2.1)
\end{tcolorbox}

\begin{tcolorbox}[resultbox]
A propriedade \'e \textbf{Limit Cycle} (ciclo-limite), definida no artigo exatamente como
oscila\c{c}\~ao persistente na sa\'ida de um sistema com entrada constante (o artigo
verifica, em particular, ciclos-limite de entrada nula --- \emph{zero-input} --- para
qualquer condi\c{c}\~ao inicial).
\end{tcolorbox}

\newpage

%==============================================================
% EXERCÍCIO 2
%==============================================================
\section{Arquitetura e completude do DSVerifier (Artigo A)}

\subsection{(a) Fluxo de verificação e posição do módulo}

\begin{tcolorbox}[quotebox]
``DSVerifier is an internal module for the efficient SMT-based context-bounded model checker
(ESBMC) [\ldots]. The complete verification tool includes four components from ESBMC [\ldots]:
C parser, GOTO Program, GOTO Symex, and SMT solver. [\ldots] The DSVerifier module is included
\emph{before} the C parser.'' (Artigo A, Se\c{c}\~ao~2, Fig.~1)
\end{tcolorbox}

Fluxo completo: \emph{Digital System Specification} $\to$ \textbf{DSVerifier}
(internamente: Initialization $\to$ Validation $\to$ Instrumentation) $\to$ arquivo ANSI-C
$\to$ \textbf{C Parser} $\to$ \textbf{GOTO Program} $\to$ \textbf{GOTO Symex} $\to$
\textbf{SMT Solver} (Boolector).

\begin{tcolorbox}[notabox]
\'E a \'unica posi\c{c}\~ao poss\'ivel porque o que o m\'odulo \textbf{produz} \'e,
literalmente, um arquivo-texto ANSI-C (``an ANSI-C code file is generated [\ldots] this file
is directly sent to the C parser module''). O pr\'oximo est\'agio capaz de consumir texto C
\'e, por defini\c{c}\~ao, o C Parser --- o \emph{front-end} de qualquer \emph{toolchain}
baseada em fonte C. N\~ao pode entrar depois do C Parser (GOTO Program j\'a espera uma
representa\c{c}\~ao intermedi\'aria, n\~ao texto-fonte), nem depois do GOTO Symex ou no SMT
Solver (o DSVerifier n\~ao tem nenhuma no\c{c}\~ao de IR-GOTO ou de f\'ormula SMT --- ele
s\'o sabe emitir C). A posi\c{c}\~ao \'e ditada inteiramente pelo tipo do artefato de sa\'ida.
\end{tcolorbox}

\subsection{(b) ANSI-C em vez de SMT direto: uma vantagem e uma limitação}

\textbf{Vantagem.} O DSVerifier reaproveita integralmente o \emph{pipeline} maduro do ESBMC
(parser, tradu\c{c}\~ao GOTO, execu\c{c}\~ao simb\'olica, codifica\c{c}\~ao SMT e m\'ultiplos
\emph{back-ends} de solver intercambi\'aveis) sem precisar implementar seu pr\'oprio
codificador SMT ou sem\^antica de \emph{bit-vectors}; basta saber gerar C corretamente. Como o
pr\'oprio artigo nota, o arquivo gerado ``can be verified by any C model-checker that supports
bit-vector reasoning'' --- portabilidade/independ\^encia de ferramenta.

\textbf{Limita\c{c}\~ao.} A qualidade da verifica\c{c}\~ao fica inteiramente refém da
sem\^antica gen\'erica de C que o ESBMC oferece: o DSVerifier n\~ao pode expressar nem
explorar nenhuma otimiza\c{c}\~ao espec\'ifica de dom\'inio (por exemplo, um procedimento de
decis\~ao dedicado a aritm\'etica de ponto fixo Q-format), pois tudo precisa antes ser
reduzido a opera\c{c}\~oes inteiras/\emph{bit-vector} de C puro; qualquer imprecis\~ao,
inconsist\^encia de solidez (como a incompletude discutida no Artigo B) ou escolha de modelo
de mem\'oria do ESBMC \'e herdada de forma transparente, sem que o DSVerifier tenha qualquer
oportunidade de intervir no n\'ivel da f\'ormula SMT.

\subsection{(c) A assimetria entre estabilidade/fase mínima e as demais propriedades}

\begin{tcolorbox}[quotebox]
``The verification of stability and minimum-phase is complete and sound, since it does not
depend on system outputs and inputs. However, the verification of other property types is
typically unsound, unless some induction technique is used.'' (Artigo A, Se\c{c}\~ao~2.2)
\end{tcolorbox}

Estabilidade e fase m\'inima s\~ao propriedades \textbf{puramente alg\'ebricas dos
coeficientes} do filtro (localiza\c{c}\~ao de p\'olos/zeros no plano $z$): dado um conjunto
\textbf{fixo e finito} de coeficientes quantizados, verificar se todos os p\'olos (ou zeros)
est\~ao dentro do c\'irculo unit\'ario \'e um c\'alculo fechado, de tamanho constante, que
\textbf{dispensa completamente o racioc\'inio sobre execu\c{c}\~oes} --- n\~ao \'e preciso
desenrolar nenhum la\c{c}o de controle nem simular a opera\c{c}\~ao do filtro ao longo do
tempo; n\~ao h\'a ``$k$'' a escolher, porque n\~ao h\'a sequ\^encia temporal de entradas/sa\'idas
envolvida.

\begin{tcolorbox}[notabox]
Overflow e ciclo-limite, ao contr\'ario, s\~ao propriedades \textbf{sobre execu\c{c}\~oes}:
perguntam se \emph{alguma} sequ\^encia de entradas, ao longo de um n\'umero potencialmente
ilimitado de amostras, algum dia causa overflow ou faz o sistema entrar em oscila\c{c}\~ao
persistente. Isso exige desenrolar a equa\c{c}\~ao de diferen\c{c}as do filtro por um
horizonte temporal que, em opera\c{c}\~ao real, \'e ilimitado --- exatamente o cen\'ario em
que BMC puro (desenrolamento at\'e $k$) \'e estruturalmente incompleto: um resultado
``sucesso'' em profundidade $k$ nunca descarta uma viola\c{c}\~ao no passo $k+1$. Da\'i a
necessidade de uma t\'ecnica de indu\c{c}\~ao (Artigo B) para elevar a garantia limitada a
$k$ passos para uma garantia v\'alida \emph{para todo tempo}.
\end{tcolorbox}

\subsection{(d) Aumentar \texttt{-{}-x-size} não estabelece corretude}

O par\^ametro \vv{--x-size} \'e o limite $k$ de desenrolamento (``the verification bound, i.e.,
the number of times the digital system will be unwound''). Um resultado \vv{successful} com
um dado \vv{--x-size=k} licencia apenas a afirma\c{c}\~ao \textbf{limitada}: ``nenhuma
viola\c{c}\~ao da propriedade foi encontrada dentro das primeiras $k$ itera\c{c}\~oes de
opera\c{c}\~ao do sistema digital'' --- o equivalente a $Base_k$ insatisfat\'ivel em BMC puro.

\begin{tcolorbox}[resultbox]
Aumentar \vv{--x-size} at\'e a verifica\c{c}\~ao passar \textbf{nunca} estabelece corretude do
projeto \emph{para todo tempo de opera\c{c}\~ao}, porque overflow/ciclo-limite n\~ao t\^em um
limiar de completude (\emph{completeness threshold}) conhecido a priori: sempre \'e poss\'ivel,
em princ\'ipio, que a viola\c{c}\~ao s\'o se manifeste em $k+1$, $k+50$ ou qualquer profundidade
maior que a testada. Um ``successful'' com \vv{--x-size=k} prova estritamente que \textbf{n\~ao
h\'a} contraexemplo nos primeiros $k$ passos --- nunca que n\~ao existe contraexemplo algum.
S\'o o passo indutivo do Artigo B (Exerc\'icio~5) fornece uma prova v\'alida para
\emph{qualquer} n\'umero de itera\c{c}\~oes.
\end{tcolorbox}

\newpage

%==============================================================
% EXERCÍCIO 3
%==============================================================
\section{Os três casos da $k$-indução (Artigo B)}

\begin{lstlisting}[style=cstyle,caption={Programa motivador, Artigo B (Fig.\ 2)}]
int main(void) {
  uint64_t i = 1, sn = 0;
  uint32_t n;
  assume(n >= 1);
  while (i <= n) {
    sn = sn + a;
    i++;
  }
  assert(sn == n * a);
}
\end{lstlisting}

\subsection{(a) Por que BMC tradicional falha e quantos desdobramentos seriam necessários}

\vv{n} \'e \vv{uint32\_t} e n\~ao recebe cota superior al\'em de \vv{assume(n>=1)}, isto \'e,
$n \in [1,\,2^{32}-1]$. BMC tradicional fixa um limite $k$ de desenrolamento; para qualquer $k$
escolhido, existe um valor de $n>k$ para o qual o la\c{c}o ainda n\~ao terminou ap\'os $k$
itera\c{c}\~oes --- a \emph{unwinding assertion} correspondente falharia, ou o BMC simplesmente
nunca cobre esse caso. Como o pr\'oprio Artigo B afirma:

\begin{tcolorbox}[quotebox]
``BMC techniques find this difficult as the number of loop iterations, represented by \texttt{n},
is non-deterministically chosen; full unwinding would require $2^{32}-1$ unfoldings.''
(Artigo B, Se\c{c}\~ao~3)
\end{tcolorbox}

\begin{tcolorbox}[resultbox]
S\~ao necess\'arios $2^{32}-1$ desdobramentos para uma verifica\c{c}\~ao exaustiva --- o
n\'umero vem diretamente do maior valor represent\'avel por \vv{uint32\_t} ($2^{32}-1$), j\'a
que o la\c{c}o executa no m\'aximo \vv{n} vezes. Nenhum $k$ fixo e finito, na pr\'atica,
alcan\c{c}a esse limite em tempo/mem\'oria vi\'aveis --- da\'i a necessidade de $k$-indu\c{c}\~ao.
\end{tcolorbox}

\subsection{(b) Transformação do caso base para $k=2$}

\begin{tcolorbox}[quotebox]
``[\ldots] the base case simply inserts an unwinding assumption to the respective loop-free
program $P'$, consisting of the termination condition $\sigma$ after the loop, to ensure that
it finds a counterexample of depth $k$ without reporting any false incorrect result (i.e., to
avoid unfolding loops partially).'' (Artigo B, Se\c{c}\~ao~4.1.2)
\end{tcolorbox}

\begin{lstlisting}[style=cstyle,caption={Caso base, $k=2$ (padr\~ao da Fig.\ 9 do Artigo B)}]
int main(void) {
  uint64_t i, sn = 0;
  uint32_t n = nondet_uint();
  assume(n >= 1);
  i = 1;
  if (i <= n) {            // copia 1
    sn = sn + a;
    i++;
    if (i <= n) {          // copia 2
      sn = sn + a;
      i++;
    }
  }
  assume(i > n);           // unwinding assumption
  assert(sn == n * a);
}
\end{lstlisting}

\begin{tcolorbox}[resultbox]
A instru\c{c}\~ao inserida ao final \'e um \textbf{\vv{assume(i > n)}} (n\~ao um \vv{assert}).
Seu papel \'e \textbf{descartar} exatamente as execu\c{c}\~oes em que o la\c{c}o \textbf{ainda
n\~ao terminou} ap\'os as $k=2$ c\'opias (isto \'e, os casos com $n>k$, onde \vv{sn} s\'o
acumulou $k\cdot a$ em vez de $n\cdot a$). Sem esse \vv{assume}, o \vv{assert(sn==n*a)}
falharia trivialmente para qualquer $n>k$ --- n\~ao por um bug real, mas simplesmente porque o
la\c{c}o foi desenrolado poucas vezes demais. O \vv{assume} \'e o que \textbf{impede que a
ferramenta reporte um contraexemplo falso} (falso positivo de bug), causado por
desenrolamento parcial, garantindo que o caso base s\'o julgue execu\c{c}\~oes genuinamente
completas ($n\le k$).
\end{tcolorbox}

\subsection{(c) Transformação da condição de avanço para $k=2$}

\begin{lstlisting}[style=cstyle,caption={Condição de avanço (\emph{forward condition}), $k=2$ (padrão da Fig.\ 10 do Artigo B)}]
int main(void) {
  uint64_t i, sn = 0;
  uint32_t n = nondet_uint();
  assume(n >= 1);
  i = 1;
  if (i <= n) {
    sn = sn + a;
    i++;
    if (i <= n) {
      sn = sn + a;
      i++;
    }
  }
  assert(i > n);           // check loop invariant
  assert(sn == n * a);
}
\end{lstlisting}

\begin{tcolorbox}[resultbox]
\textbf{N\~ao \'e a mesma instru\c{c}\~ao.} O caso base insere um \vv{assume(i>n)}; a
condi\c{c}\~ao de avan\c{c}o insere um \textbf{\vv{assert(i>n)}} --- ``\emph{check loop
invariant}'', nas palavras do pr\'oprio c\'odigo-figura do artigo. A diferen\c{c}a de
significado \'e o ato de fala: o \vv{assume} do caso base \emph{descarta} (filtra) as
execu\c{c}\~oes n\~ao-conclu\'idas para evitar falso positivo de bug; o \vv{assert} da
condi\c{c}\~ao de avan\c{c}o \emph{prova ativamente} que o la\c{c}o \textbf{sempre} termina
dentro de $k$ itera\c{c}\~oes (isto \'e, tenta demonstrar que $n\le k$ \'e uma consequ\^encia
necess\'aria, o que s\'o \'e poss\'ivel se o la\c{c}o for de fato limitado por $k$). Para este
programa espec\'ifico, como \vv{n} \'e irrestrito em $[1,2^{32}-1]$, a condi\c{c}\~ao de
avan\c{c}o \textbf{nunca} ter\'a sucesso para nenhum $k$ finito (sempre existe $n=k+1$) --- \'e
precisamente por isso que a prova de corretude deste programa depende do \textbf{passo
indutivo} (Exerc\'icio~4), n\~ao da condi\c{c}\~ao de avan\c{c}o.
\end{tcolorbox}

\subsection{(d) Por que a ordem base $\to$ avanço/indutivo importa}

\begin{tcolorbox}[quotebox]
``The algorithm runs up to a maximum number of iterations and only increases the value of $k$
if it can not falsify the property during the base case. Note that $k$ is incremented only at
the start of the \texttt{else} branch in line~9 of Fig.~3.'' (Artigo B, Se\c{c}\~ao~4.1)
\end{tcolorbox}

Pelo pseudoc\'odigo da Fig.~3, a cada itera\c{c}\~ao do algoritmo: primeiro roda-se
\vv{base\_case(P, phi, k)}; se ele encontra um contraexemplo, o algoritmo retorna
\vv{false} \textbf{imediatamente}; s\'o no ramo \vv{else} (caso base \emph{n\~ao} achou bug)
\'e que \vv{k} \'e incrementado e ent\~ao se tenta \vv{forward\_condition} e, falhando essa,
\vv{inductive\_step}.

\begin{tcolorbox}[notabox]
A ordem importa por \textbf{corretude}, n\~ao s\'o por efici\^encia. A validade cl\'assica da
$k$-indu\c{c}\~ao (Eq.~2 do Artigo B, Se\c{c}\~ao~2.3) exige a conjun\c{c}\~ao de \textbf{duas}
condi\c{c}\~oes: $Base_k$ insatisfat\'ivel (nenhuma viola\c{c}\~ao alcan\c{c}\'avel do estado
inicial em at\'e $k$ passos) \textbf{e} $Step_k$ insatisfat\'ivel (nenhuma transi\c{c}\~ao de
$k$ estados bons para um estado ruim). O algoritmo da Fig.~3 obt\'em a metade ``$Base_k$'' de
forma \emph{acumulada}: a cada itera\c{c}\~ao do la\c{c}o externo, o caso base j\'a foi
checado (sem bug) para todos os valores de $k$ percorridos at\'e ali. Se a condi\c{c}\~ao de
avan\c{c}o ou o passo indutivo fossem executados \textbf{primeiro}, o algoritmo poderia
retornar \vv{true} baseado apenas em $Step_k$ insatisfat\'ivel, \textbf{sem} nunca ter
confirmado $Base_k$ insatisfat\'ivel no mesmo $k$ --- ou seja, poderia reportar corretude
mesmo havendo um bug genu\'ino, alcan\c{c}\'avel do estado inicial em poucos passos, que
simplesmente nunca chegou a ser checado. Checar o caso base primeiro n\~ao \'e apenas mais
barato computacionalmente (o pr\'oprio artigo nota que rodar a condi\c{c}\~ao de avan\c{c}o ou
o passo indutivo com $k$ pequeno demais desperdi\c{c}a recursos, j\'a que ``loops are usually
unfolded at least two times'') --- \'e o que garante que todo veredito \vv{true} tem, de fato,
tanto a \^ancora ($Base_k$) quanto o passo ($Step_k$) da indu\c{c}\~ao matem\'atica estabelecidos.
\end{tcolorbox}

\newpage

%==============================================================
% EXERCÍCIO 4
%==============================================================
\section{O passo indutivo e suas transformações (Artigo B)}

O la\c{c}o \vv{while(c) \{E;\}} \'e convertido em \vv{A; while(c) \{ S; E; U; R; \}} (Eq.~4).

\subsection{(a) Função de cada bloco}

\begin{tcolorbox}[quotebox]
``$A$ is the code (or sequence of commands) responsible for assigning non-deterministic values
to all loops variables, i.e., the state is havocked before the loop, $c$ is the termination
condition of the loop \texttt{while}, $S$ is the code to store the current state of the
program variables before executing the statements of $E$, $E$ is the original body of the
\texttt{while} loop, $U$ is the code to update all state variables with local values after
executing $E$, and $R$ is the code to remove redundant states.'' (Artigo B, Se\c{c}\~ao~4.1.2)
\end{tcolorbox}

\begin{center}
\begin{tabular}{@{}p{1.3cm}p{12.8cm}@{}}
\toprule
Bloco & Função \\
\midrule
$A$ & \emph{Havoc}: atribui valores n\~ao determin\'isticos a todas as vari\'aveis do la\c{c}o, \textbf{antes} do la\c{c}o começar --- generaliza o estado inicial para que o model checker explore implicitamente todas as configura\c{c}\~oes poss\'iveis. \\
$S$ & Salva (armazena) o estado corrente das vari\'aveis do programa \textbf{antes} de executar $E$ na itera\c{c}\~ao atual --- popula o vetor de estados \vv{sv}. \\
$U$ & Atualiza todas as vari\'aveis de estado com os valores locais \textbf{depois} de executar $E$. \\
$R$ & Remove estados redundantes (repetidos), via um \vv{assume} que compara o vetor de estados salvo com o estado atual. \\
\bottomrule
\end{tabular}
\end{center}

\subsection{(b) O bloco $R$ e a terminação do passo indutivo}

O bloco $R$ \'e um \vv{assume} do tipo $sv_1[0]\neq cs_1 \wedge sv_1[0]\neq cs_1\wedge
sv_2[1]\neq cs_2 \wedge \ldots$ (Eq.~8), isto \'e, exige que o estado da itera\c{c}\~ao atual
seja \textbf{diferente} de todos os estados salvos nas itera\c{c}\~oes anteriores.

\begin{tcolorbox}[notabox]
$R$ elimina execu\c{c}\~oes em que a sequ\^encia de $k$ estados percorridos \textbf{n\~ao \'e
um caminho simples} (isto \'e, revisita um estado j\'a visto --- um ciclo dentro da janela de
$k$ itera\c{c}\~oes). Isso \'e necess\'ario porque a validade cl\'assica da $k$-indu\c{c}\~ao
sobre sistemas de transi\c{c}\~ao (Sheeran et al.\ 2000; E\'en \& S\"orensson 2003, citados no
Artigo B) exige que os $k{+}1$ estados testados no passo indutivo sejam
\textbf{genuinamente distintos}: sem essa restri\c{c}\~ao, o SMT solver poderia ``provar'' o
passo indutivo trivialmente sobre um \textbf{ciclo curto} que se repete indefinidamente sem
progresso real, produzindo uma conclus\~ao \emph{unsound} (o passo pareceria v\'alido para
qualquer $n$ apenas porque o mesmo pequeno conjunto de estados foi revisitado, sem cobrir o
comportamento real do programa para valores maiores). Ao proibir estados repetidos, $R$
garante que, quando o passo indutivo termina com sucesso, a resposta \'e de fato \'util
(sound) --- corresponde a uma cadeia real e crescente de estados do programa, n\~ao a um
artefato de um la\c{c}o de estados degenerado.
\end{tcolorbox}

\subsection{(c) Havoc apenas das variáveis do laço vs.\ Große et al.\ (todas as variáveis)}

\begin{tcolorbox}[quotebox]
``Similar to Donaldson et al.\ [8,9], we havoc only the variables that occur in the loop, i.e.,
all loops variables are assigned non-deterministic values. [\ldots] Differently from our
approach, Gro\ss{}e et al.\ havoc all program variables, which makes it difficult to check for
the reachability of an error since they do not provide enough information to constrain the
havocked variables in the program.'' (Artigo B, Se\c{c}\~ao~4)
\end{tcolorbox}

\begin{tcolorbox}[resultbox]
Ao randomizar \textbf{todas} as vari\'aveis do programa (inclusive as que o la\c{c}o nunca
toca), o m\'etodo de Gro\ss{}e et al.\ destr\'oi qualquer rela\c{c}\~ao/invariante que essas
vari\'aveis n\~ao-relacionadas ao la\c{c}o mantinham com o restante do estado --- n\~ao
sobra informa\c{c}\~ao suficiente para \textbf{restringir} os valores havockeados de forma
coerente com o que \'e de fato alcan\c{c}\'avel. Na pr\'atica, isso torna imposs\'ivel
raciocinar com precis\~ao sobre se um estado de erro \'e \textbf{realmente} alcan\c{c}\'avel:
o solver passa a considerar combina\c{c}\~oes de valores para vari\'aveis n\~ao-relacionadas
ao la\c{c}o que nunca ocorreriam em uma execu\c{c}\~ao real, inflando artificialmente o
espa\c{c}o de estados considerado (com risco de falsos alarmes) ou impedindo a prova de
alcan\c{c}abilidade (por faltar contexto). Havockear apenas as \emph{vari\'aveis do la\c{c}o}
(Defini\c{c}\~ao~2/3 do Artigo B) preserva intactas as rela\c{c}\~oes entre o restante do
estado do programa e a l\'ogica externa ao la\c{c}o, mantendo o racioc\'inio sobre
alcan\c{c}abilidade preciso.
\end{tcolorbox}

\subsection{(d) Por que a hipótese $\{Post \wedge \neg(i \le n)\}$ prova o caso geral}

\begin{tcolorbox}[quotebox]
``Lemma 1 If the induction hypothesis $\{Post\wedge\neg(i\le n)\}$ holds for $k+1$ consecutive
iterations, then it also holds for $k$ preceding iterations. [\ldots] the SMT solver can
easily verify Lemma~1 and conclude that \texttt{sn == n*a} is an inductive invariant relative
to \texttt{n}.'' (Artigo B, Se\c{c}\~ao~4.3.3)
\end{tcolorbox}

A chave \'e que, no passo indutivo, o bloco $A$ atribui valores \textbf{n\~ao-determin\'isticos
gen\'ericos} \`as vari\'aveis do la\c{c}o (via \vv{nondet\_uint()}) --- o racioc\'inio n\~ao
fixa nenhum valor concreto de \vv{n}, nem assume que o estado inicial da janela de $k+1$
itera\c{c}\~oes \'e alcan\c{c}ado a partir de um \vv{n} espec\'ifico. O que \'e provado \'e um
\textbf{esquema geral}: ``\emph{para quaisquer} $k{+}1$ estados consecutivos que satisfa\c{c}am
$Post\wedge\neg(i\le n)$, o mesmo vale para os $k$ estados anteriores''.

\begin{tcolorbox}[resultbox]
Combinando essa hip\'otese gen\'erica (v\'alida para qualquer janela de $k{+}1$ itera\c{c}\~oes,
n\~ao amarrada a um \vv{n} concreto) com o caso base --- que j\'a estabeleceu
\vv{sn==n*a} para todo $n\le k$ --- obt\'em-se, por indu\c{c}\~ao matem\'atica cl\'assica
(base + passo), que a propriedade vale para \textbf{qualquer} $n\ge1$, e n\~ao apenas at\'e o
limite $k$ explorado explicitamente. Isso \'e exatamente o que separa $k$-indu\c{c}\~ao de BMC
puro: o passo indutivo, uma vez estabelecido, generaliza para infinitos valores de \vv{n} sem
precisar desenrolar mais nenhum deles.
\end{tcolorbox}

\newpage

%==============================================================
% EXERCÍCIO 5
%==============================================================
\section{Integração entre os dois artigos}

\subsection{(a) Qual dos três casos fecha a lacuna de solidez do Artigo A}

\'E o \textbf{passo indutivo}. O caso base apenas encontra bugs at\'e profundidade $k$ ---
n\~ao contribui nenhuma garantia positiva de aus\^encia de erro al\'em de $k$. A
\textbf{condi\c{c}\~ao de avan\c{c}o} prova a propriedade ``em todos os estados alcan\c{c}\'aveis
dentro de $k$ itera\c{c}\~oes'' --- ainda uma garantia \textbf{limitada} a $k$, id\^entica em
esp\'irito ao BMC puro que o Artigo A j\'a descreve como insuficiente (``t\'ipicamente n\~ao
s\'olida''). Apenas o \textbf{passo indutivo}, ao provar que a validade em $k$ itera\c{c}\~oes
consecutivas implica validade em $k{+}1$ (Lema~1), e combinado com o caso base como \^ancora,
produz uma conclus\~ao v\'alida para \textbf{qualquer} n\'umero de itera\c{c}\~oes --- a
``t\'ecnica de indu\c{c}\~ao'' que o Artigo A antecipa como necess\'aria para tornar
overflow/ciclo-limite tão s\'olidos quanto estabilidade/fase m\'inima.

\begin{tcolorbox}[notabox]
\textbf{Os outros dois casos s\~ao dispens\'aveis \emph{para esse fim espec\'ifico}?}
A condi\c{c}\~ao de avan\c{c}o \'e dispens\'avel para fechar a lacuna de solidez em sistemas
genuinamente de opera\c{c}\~ao ilimitada (ela s\'o ajuda quando o la\c{c}o de fato termina
dentro de algum $k$ pequeno, um atalho de efici\^encia, n\~ao uma fonte de solidez
\emph{unbounded}). J\'a o \textbf{caso base n\~ao \'e dispens\'avel}, mesmo para esse fim
espec\'ifico: pela Eq.~2 do Artigo B, a validade da $k$-indu\c{c}\~ao exige $Base_k$ insatisfat\'ivel
\textbf{junto com} $Step_k$ insatisfat\'ivel --- sem confirmar o caso base, o passo indutivo
sozinho n\~ao garante aus\^encia de bug nos primeiros $k$ passos (ver Exerc\'icio~3d). O caso
base \'e a \^ancora l\'ogica sem a qual o passo indutivo n\~ao licencia nenhuma conclus\~ao
sobre o programa completo.
\end{tcolorbox}

\subsection{(b) Alterações necessárias no fluxo do DSVerifier}

Pela Fig.~1 do Artigo A, a $k$-indu\c{c}\~ao do Artigo B \'e implementada como transforma\c{c}\~oes
aplicadas \textbf{no n\'ivel da representa\c{c}\~ao intermedi\'aria GOTO-program}, durante a
convers\~ao C/C++ $\to$ GOTO (``these transformations take place at the intermediate
representation level, during the conversion of the C/C++ program into a GOTO-program''), ou
seja, \textbf{dentro} do pr\'oprio ESBMC, entre os est\'agios ``C Parser'' e ``GOTO Symex''.

\begin{tcolorbox}[resultbox]
O m\'odulo DSVerifier em si (Initialization $\to$ Validation $\to$ Instrumentation, antes do
C Parser) \textbf{n\~ao precisa de nenhuma altera\c{c}\~ao}: o arquivo ANSI-C que ele gera
continua id\^entico. A mudan\c{c}a ocorre exclusivamente na \textbf{etapa 5} da metodologia da
Fig.~2 do Artigo A (``Verify Using a BMC tool''): em vez de invocar o ESBMC em modo de BMC puro
(desenrolamento simples at\'e \vv{--x-size}), invoca-se o ESBMC com a $k$-indu\c{c}\~ao
habilitada (na pr\'atica, as flags \vv{--k-induction --k-step <k>} usadas nos experimentos do
Artigo B). Como as transforma\c{c}\~oes de $k$-indu\c{c}\~ao agem automaticamente sobre
\emph{qualquer} la\c{c}o presente no GOTO-program, elas se aplicam tamb\'em aos la\c{c}os de
verifica\c{c}\~ao que a pr\'opria instrumenta\c{c}\~ao do DSVerifier injeta, sem que o
DSVerifier precise ``saber'' que a $k$-indu\c{c}\~ao existe.
\end{tcolorbox}

\subsection{(c) O código gerado pelo DSVerifier respeita as restrições do passo indutivo?}

\begin{tcolorbox}[quotebox]
``[\ldots] the inductive step of our $k$-induction algorithm does not support heap-manipulating
programs and unbounded recursion [\ldots] Most of the incorrect results provided by ESBMC are
related to bugs in the coding of the inductive step transformations, because we do not support
\texttt{continue} and \texttt{return} (from inside a loop) statements.'' (Artigo B, Introdu\c{c}\~ao
e Se\c{c}\~ao~5.2)
\end{tcolorbox}

\begin{tcolorbox}[resultbox]
\textbf{Sim, respeita.} As realiza\c{c}\~oes de filtros/controladores digitais que o
DSVerifier gera (formas diretas, formas delta e cascata --- vetores de coeficientes
\vv{ds.a}/\vv{ds.b} de tamanho fixo conhecido em tempo de gera\c{c}\~ao, via \vv{a\_size}/
\vv{b\_size}) s\~ao n\'ucleos num\'ericos \textbf{est\'aticos, simples e n\~ao-recursivos}: um
la\c{c}o (ou sequ\^encia de la\c{c}os) que acumula somas e produtos sobre \emph{arrays} de
tamanho fixo, sem aloca\c{c}\~ao din\^amica (\emph{heap}), sem chamadas recursivas (a
realiza\c{c}\~ao de um filtro \'e inerentemente iterativa, espelhando diretamente a
equa\c{c}\~ao de diferen\c{c}as) e sem desvios de controle antecipados (\vv{continue}/\vv{return})
dentro do la\c{c}o de processamento de amostras.
\end{tcolorbox}

A caracter\'istica decisiva \'e que c\'odigo de realiza\c{c}\~ao de sistemas digitais \'e um
\textbf{``la\c{c}o de aritm\'etica em linha reta''}: em cada itera\c{c}\~ao executa-se sempre a
mesma sequ\^encia fixa de multiplica\c{c}\~oes e somas sobre estados/coeficientes de tamanho
est\'atico, sem ramifica\c{c}\~ao de controle que dependa de dados nem estrutura de dados
din\^amica --- exatamente a classe de programas para a qual as transforma\c{c}\~oes do
Artigo B (havoc de vari\'aveis escalares/vetoriais de tamanho fixo conhecido, bloco $R$ sobre
um \vv{statet} tamb\'em de tamanho fixo) foram desenhadas.

\subsection{(d) Mais resultados corretos \emph{e} mais \emph{wrong proofs}: o que pesa mais?}

\begin{tcolorbox}[quotebox]
``Wrong proofs is the number of false positive results (i.e., the tool reports SAFE
incorrectly).'' (Artigo B, Se\c{c}\~ao~5.2) --- Tabela~1: ESBMC paralelo, 3153 resultados
corretos / 265 \emph{wrong proofs}; CPAchecker, 2349 resultados corretos / 26 \emph{wrong
proofs}.
\end{tcolorbox}

\begin{center}
\begin{tabular}{@{}lccc@{}}
\toprule
& ESBMC sequencial & ESBMC paralelo & CPAchecker \\
\midrule
Resultados corretos & 3071 & \textbf{3153} & 2349 \\
\emph{Wrong proofs} (falso SAFE) & 111 & 265 & \textbf{26} \\
\bottomrule
\end{tabular}
\end{center}

\begin{tcolorbox}[resultbox]
Para um engenheiro decidindo se adota a ferramenta na verifica\c{c}\~ao de um controlador
embarcado \textbf{cr\'itico}, o n\'umero de \emph{wrong proofs} pesa muito mais do que o
n\'umero bruto de resultados corretos --- e a raz\~ao \'e a assimetria de consequ\^encias
entre os dois tipos de erro:
\begin{itemize}[leftmargin=1.6cm]
  \item Um \textbf{timeout/unknown} (nenhum veredito) \'e um modo de falha \emph{seguro}: a
  engenharia simplesmente n\~ao ganha um certificado formal e recorre a outra t\'ecnica
  (teste, revis\~ao manual, outra ferramenta) --- nenhum dano \'e feito, nenhuma confian\c{c}a
  falsa \'e criada.
  \item Uma \textbf{\emph{wrong proof}} \'e um modo de falha \emph{perigoso}: a ferramenta
  certifica como SAFE um controlador que na verdade cont\'em um defeito. Isso gera
  \textbf{confian\c{c}a falsa} --- exatamente o cen\'ario que a verifica\c{c}\~ao formal existe
  para eliminar. Um controlador defeituoso pode ser embarcado em produ\c{c}\~ao \emph{acreditando-se}
  que foi formalmente verificado.
\end{itemize}
Relativamente, o ESBMC paralelo erra com falso-SAFE em $265/3153\approx8{,}4\%$ dos casos
resolvidos, contra $26/2349\approx1{,}1\%$ do CPAchecker --- quase $8\times$ mais frequente.
Vale notar ainda que a pr\'opria implementa\c{c}\~ao paralela do ESBMC \textbf{piora} esse
n\'umero em rela\c{c}\~ao \`a sequencial (111$\to$265 \emph{wrong proofs}), precisamente porque
processos que estourariam mem\'oria/tempo na vers\~ao sequencial (e ficariam honestamente
``unknown'') s\~ao convertidos em vereditos definitivos na vers\~ao paralela --- um ganho de
cobertura obtido \`a custa de solidez. Para certifica\c{c}\~ao de um sistema cr\'itico, a
escolha razo\'avel \'e privilegiar a menor taxa de falso-SAFE (CPAchecker, ou ESBMC sequencial
em vez do paralelo), mesmo abrindo m\~ao de cobertura bruta --- e, idealmente, usar
m\'ultiplas ferramentas/t\'ecnicas em conjunto, j\'a que nenhuma delas tem taxa zero de
\emph{wrong proofs}.
\end{tcolorbox}

\newpage

%==============================================================
% REFERÊNCIAS
%==============================================================
\begin{thebibliography}{9}
\bibitem{dsverifier} H.~I.~Ismail, I.~V.~Bessa, L.~C.~Cordeiro, E.~B.~de Lima Filho,
  J.~E.~Chaves Filho. \emph{DSVerifier: A Bounded Model Checking Tool for Digital Systems}.
  SPIN 2015, LNCS 9232, pp.~126--131.
\bibitem{kinduction} M.~Y.~R.~Gadelha, H.~I.~Ismail, L.~C.~Cordeiro. \emph{Handling Loops in
  Bounded Model Checking of C Programs via $k$-Induction}. STTT, 2017.
\end{thebibliography}

\end{document}
